% Αγγλική Περίληψη

\begin{abstracteng}\setlanguage{american}

Choosing the appropriate location is a critical success factor for every business activity. The purpose of this thesis is to develop a system that combines the power of spatial data with the capabilities of genetic artificial intelligence, with the aim of supporting the optimal choice of business location. The proposed system utilizes geospatial characteristics, such as population density, accessibility, as well as demographic or economic data of areas. Through a combination of spatial analysis and advanced language models (LLMs), the system extracts information from heterogeneous sources, processes it, and produces explanatory and personalized location suggestions for new businesses. The study includes the construction of two pipelines based on LLMs, the selection of appropriate geospatial features, the creation of a search index, the training of a location evaluation model, and the final generation of responses. At the same time, the effectiveness of the system is evaluated through accuracy metrics.

\vspace*{\fill}
\noindent{\large\bf{Keywords:}}\\ 
term, term, $\dots$, term
\end{abstracteng}